\chapter{Experiments and Discussion}
\label{cha:experiments_discussion}


\section{Experimental Setup}
\label{sec:experimental_setup}

To rigorously evaluate the proposed multi-year architecture against single-year baselines, a standardized experimental framework was established.
All models were trained, validated, and tested using identical data splits and evaluation metrics to ensure a fair comparative analysis.

\subsection{HPC Environment and I/O Bottleneck Mitigation}
\label{subsec:hpc_environment}

Training deep sequence models on massive, high-dimensional satellite time series is computationally demanding.
To facilitate this, all experiments were executed on a dedicated High-Performance Computing (HPC) node equipped with an AMD Ryzen 7 5800X processor,
67GB of system RAM, and an NVIDIA GeForce RTX 3060 GPU, running on an Ubuntu 22.04 LTS environment.

A primary challenge in processing the INVEKOS dataset is the severe input/output (I/O) bottleneck caused by loading thousands of multi-year temporal sequences.
Initial exploratory experiments attempted to utilize the full 2D spatial patches of the agricultural parcels.
However, loading continuous multi-year, multi-spectral spatial patches introduced an insurmountable I/O bottleneck.
Because the 2D patch dataset was exceptionally large (approximately 1.3 TB), it could not be stored on the
HPC node's local solid-state drive (SSD) and instead had to be hosted on a shared network folder.
The immense volume of network-based disk reads required to process the 2D geometries caused severe GPU starvation,
rendering the training process computationally infeasible despite the HPC infrastructure.

To resolve this, the experimental framework was adapted to utilize the highly optimized 1D pixel LMDB dataset provided by Perantoni et al. \cite{perantoni_bayesian_2024}.
As detailed in the methodology, adopting this memory-mapped structure effectively eliminated the I/O bottleneck.
It allowed the PyTorch Lightning data loaders to rapidly access sequences directly from disk, sustaining high GPU utilization without the latency overhead of the 2D spatial patches.
Furthermore, all model executions, hyperparameter tuning trials via Ray Tune, and resulting performance metrics were systematically tracked and versioned using a centralized MLflow server.

\subsection{Data Loading Strategies}
\label{subsec:data_loading}

To properly evaluate the impact of multi-year memory, the experimental framework utilized two distinct data loading
strategies tailored to the architectural needs of the models.

For the single-year baseline models (1D-CNN and Bidirectional GRU), the dataset was structured using an intra-year loader.
This loader effectively decoupled the temporal sequence of the dataset, extracting and serving only a single growing season (e.g., 2018) at a time.
This constrained the baselines to rely entirely on localized, intra-year phenological signatures for classification.

In contrast, the proposed Hybrid Transformer-LSTM required a full sequence view of the crop rotation cycle.
To facilitate this, an inter-year data loader was engineered to serve the entire continuous multi-year sequence for a given agricultural parcel simultaneously.
When passed to the Hybrid model, this multi-year tensor is first processed by the Intra-Year Transformer,
which concurrently analyzes the phenology within each individual year.
The summarized yearly representations are then passed to the Inter-Year LSTM to sequentially process the multi-year crop rotation.
This explicit separation in data handling guarantees that the performance delta observed between the architectures is strictly
a product of the Hybrid model's multi-year memory capabilities.

\subsection{Evaluation Metrics}
\label{subsec:evaluation_metrics}

Agricultural land cover datasets are inherently unbalanced; a few dominant crop types (such as maize or winter wheat) often account for the vast majority of cultivated land,
while specialized crops represent a tiny fraction of the dataset.
Consequently, relying solely on overall accuracy can yield highly misleading evaluations.
To provide a comprehensive assessment, the models were evaluated using two primary metrics: \textbf{Overall Accuracy} and the \textbf{Macro F1-Score}.

\textbf{Overall Accuracy (OA)} measures the total proportion of correctly classified agricultural parcels across all available classes. It is defined as:
$$ \text{OA} = \frac{\sum_{i=1}^{C} TP_i}{N} $$
where $C$ is the total number of target classes (47 in this study), $TP_i$ represents the True Positives for class $i$, and $N$ is the total number of samples in the dataset.

\textbf{Macro F1-Score} is the primary metric used to evaluate true model performance in this study. The F1-score for a single class is the harmonic mean of Precision and Recall.
To compute the Macro F1-score, the F1-score is calculated independently for each of the 47 classes and then unweightedly averaged:
$$ \text{Macro F1} = \frac{1}{C} \sum_{i=1}^{C} \left( 2 \times \frac{\text{Precision}_i \times \text{Recall}_i}{\text{Precision}_i + \text{Recall}_i} \right) $$
Unlike a weighted F1-score, the Macro F1-score treats every class equally, regardless of its support size.
A model can only achieve a high Macro F1-score if it successfully learns to identify both the dominant and the minority crop classes,
making it the most rigorous benchmark for this highly imbalanced dataset.


\section{Baseline Results (Non-Hybrid Models)}
\label{sec:baseline_results}

To quantify the value of modeling multi-year crop rotations, the performance of the proposed architecture must first be contextualized against established single-year models.
To ensure statistical robustness and account for weight initialization variances, all baseline models were trained and evaluated across three distinct seeds.

\subsection{Temporal CNN Performance}
\label{subsec:cnn_performance}

Following the hyperparameter tuning phase, the optimal Temporal 1D-CNN was configured with a batch size of 16, a learning rate of $2.66\mathrm{e}{-3}$,
and a weight decay of $1.65\mathrm{e}{-6}$. The model was trained for 40 epochs.

As shown in Table \ref{tab:hybrid_comparison}, the CNN achieved a mean overall test accuracy of 55.17 $\pm$0.00\%.
However, the model struggled significantly with the class imbalance of the INVEKOS dataset, recording a mean test Macro F1-score of only 26.78 $\pm$0.00\%.
The large discrepancy between the overall accuracy and the F1-score indicates that while the CNN can adequately identify majority crop classes based on local temporal convolutions,
it fails to capture the defining phenological features of minority classes.

\subsection{Temporal GRU RNN Performance}
\label{subsec:gru_performance}

The Bidirectional GRU baseline, designed to explicitly model sequential dependencies, was trained under similar conditions.
The optimal configuration utilized a hidden dimension of 128 and a network depth of 3 stacked layers.
The network was trained with a learning rate of $4.73\mathrm{e}{-4}$, a weight decay of $4.98\mathrm{e}{-6}$, and a batch size of 16.

Across the three random seeds, the GRU achieved a mean overall test accuracy of 57.78 $\pm$0.00\% and a mean Macro F1-score of 34.95 $\pm$0.00\%.
Compared to the local sliding window of the CNN, the recurrent memory of the GRU provided a measurable improvement in distinguishing the complex,
irregular time gaps inherent in the Sentinel-2 acquisitions, raising the F1-score by approximately 8.17\%.
Ultimately, however, both single-year baselines demonstrate a severe performance ceiling when isolated from historical crop rotation data.


\section{Proposed Architecture Results (Hybrid Model)}
\label{sec:hybrid_results}

\subsection{Hierarchical Transformer-LSTM Performance}
\label{subsec:transformer_lstm_performance}

To evaluate the core hypothesis of this thesis---that integrating inter-year crop rotation memory improves classification over
isolated single-year models---the Hierarchical Transformer-LSTM was evaluated under an identical experimental setup. 
Following the Bayesian hyperparameter search, the optimal multi-year architecture was configured with a
Transformer embedding dimension ($d\_model$) of 256, utilizing 2 encoder layers and 4 attention heads.
The inter-year sequence was processed by a single-layer Bidirectional LSTM with a hidden dimension of 128.
The model was trained using a batch size of 16, a learning rate of $1.57\mathrm{e}{-4}$, a weight decay of $2.91\mathrm{e}{-5}$, and a dropout rate of 10\% to prevent overfitting.

Consistent with the baselines, the optimal configuration was trained across three distinct random seeds.
As detailed in Table \ref{tab:hybrid_comparison}, the hybrid architecture achieved a remarkable mean overall test accuracy of 80.65 $\pm$0.51\%.
More importantly, it achieved a mean Macro F1-score of 61.17 $\pm$0.50\%, vastly outperforming both single-year baselines. 

This dramatic F1-score improvement of over 26\% compared to the GRU confirms that multi-year memory enables the network to correctly identify
minority crop classes that were previously completely misclassified by the CNN and GRU. By understanding what was planted in the same field in previous years,
the LSTM component effectively narrows down the probability distribution of current-year crops, breaking through the performance ceiling of intra-year models.

\begin{table}[htpb]
    \centering
    \renewcommand{\arraystretch}{1.2}
    \begin{tabular}{lccc}
        \toprule
        \textbf{Model} & \textbf{Test Accuracy (\%)} & \textbf{Test Macro F1 (\%)} & \textbf{Test Loss} \\
        \midrule
        Temporal 1D-CNN (Single-Year) & 55.17 $\pm$ 0.00 & 26.78 $\pm$ 0.00 & 1.46 $\pm$ 0.00 \\
        Bidirectional GRU (Single-Year) & 57.78 $\pm$ 0.00 & 34.95 $\pm$ 0.00 & 1.40 $\pm$ 0.00 \\
        \textbf{Transformer-LSTM (Multi-Year)} & \textbf{80.65 $\pm$ 0.51} & \textbf{61.17 $\pm$ 0.50} & \textbf{0.61 $\pm$ 0.01} \\
        \bottomrule
    \end{tabular}
    \caption{Final performance comparison across all evaluated architectures. Results are averaged across 3 random seeds, displaying the mean and standard deviation.}
    \label{tab:hybrid_comparison}
\end{table}


\section{Discussion}
\label{sec:discussion}

The experimental results validate the central premise of this research: relying solely on intra-year phenology severely limits crop
classification performance, particularly on highly imbalanced datasets.
The hierarchical multi-year architecture not only surpassed the performance ceiling of the baselines but also demonstrated unexpected computational efficiencies.

\subsection{Single-Year vs. Multi-Year Capabilities}
\label{subsec:single_vs_multi_year}

The Hybrid Transformer-LSTM significantly outperforms the single-year baselines.
While the Bidirectional GRU is theoretically better suited to capture complex temporal phenology than the CNN, it only achieved a Macro F1-score of 34.95\%.
In contrast, the Hybrid model reached a mean F1-score of 61.17\%.
This massive improvement highlights a fundamental limitation of single-year models: they evaluate each growing season in a vacuum.
In agricultural reality, planting decisions are heavily governed by strict multi-year crop rotation practices designed
to prevent soil depletion (e.g., alternating nitrogen-fixing legumes with demanding crops like maize).

By feeding a sequence of yearly embeddings into the inter-year LSTM, the Hybrid model learns these transitional probabilities.
If the network is uncertain about a crop's spectral signature in year $t$, it can infer the likely crop by analyzing the confident predictions from years $t-1$ and $t+1$. 
Furthermore, despite its architectural complexity, the Hybrid model proved to be more computationally efficient during training,
completing its required epochs across all three seeds faster than the single-year models.

This efficiency can likely be attributed to the model's hierarchical design.
As originally established by Vaswani et al. \cite{attention_is_all_you_need},
self-attention mechanisms allow for significantly more parallelization than recurrent networks, which are bottlenecked by step-by-step sequential processing.
While the baseline GRU must process hundreds of individual satellite acquisitions sequentially,
the Hybrid model leverages the Intra-Year Transformer to evaluate the intra-year sequences concurrently.
By pooling these outputs into dense, single-year vectors, the sequential computation burden on the inter-year LSTM is reduced
to a sequence length of merely six steps (representing the six years of data).

\subsection{The Impact of Class Imbalance}
\label{subsec:class_imbalance}

The INVEKOS dataset, like most real-world agricultural registries, suffers from severe long-tail class imbalance,
dominated by staple crops such as winter wheat and maize \cite{perantoni_bayesian_2024}.
This imbalance is fully exposed by the baseline models.
The 1D-CNN achieved an overall accuracy of 55.17\%, but an F1-score of only 26.78\%.
This nearly 30\% discrepancy indicates a "lazy" network: as is common in highly imbalanced datasets,
the model minimizes its training loss by over-predicting the majority classes and effectively ignoring the minority crops \cite{johnson2019survey}.

The Hybrid model significantly mitigates this failure mode. While its overall accuracy rose to 80.65\%, its F1-score jumped disproportionately higher to 61.17\%.
This proves that the multi-year memory explicitly helps rescue minority classes.
Rare or specialized crops are often planted as transitional phases within a longer crop rotation cycle to manage soil health \cite{crop_rotation, osman_crop_2013}.
By recognizing the complete multi-year cycle, the LSTM is forced to predict the correct minority class to satisfy the learned rotational sequence,
preventing the model from defaulting to the statistically safe, majority-class prediction.

Ultimately, these findings demonstrate that treating crop mapping as an isolated, single-year computer vision problem discards critical agronomic context.
Fusing intra-year attention with inter-year recurrence is essential for accurate, scalable, and computationally viable agricultural monitoring.